% Nguồn SGK Toán 9 Tập hai — Bài 28, trang 72--77:
% https://cdn3.olm.vn/upload/thuvienso/4491669493-page-72-1771844558358.png
% https://cdn3.olm.vn/upload/thuvienso/4491669493-page-73-1771844558690.png
% https://cdn3.olm.vn/upload/thuvienso/4491669493-page-74-1771844559030.png
% https://cdn3.olm.vn/upload/thuvienso/4491669493-page-75-1771844559354.png
% https://cdn3.olm.vn/upload/thuvienso/4491669493-page-76-1771844559669.png
% https://cdn3.olm.vn/upload/thuvienso/4491669493-page-77-1771844560019.png
% Authority: https://olm.vn/thu-vien-so/doc-sach/shs-toan-tap-2.4491669493
%
% Nguồn SBT Toán 9 Tập hai — Bài 28, PDF trang 53--54:
% SBT TOÁN 9 TẬP 2.pdf
% SHA-256: 5ff01fd213d1adc75f9c54b4408d6a9642a4d8038e685040efef3311a196196c

\section{Đường tròn ngoại tiếp và đường tròn nội tiếp của một tam giác}

\begin{lesson}
    \subsection{Kiến thức trọng tâm}

    \begin{center}
    \begin{tblr}{
        width=\linewidth,
        colspec={X[2,l] X[5,l]},
        cells={valign=t},
        row{1}={font=\bfseries, halign=c}
    }
        Khái niệm, thuật ngữ & Kiến thức, kĩ năng \\
        \begin{minipage}[t]{\linewidth}
        \begin{itemize}
            \item Đường tròn ngoại tiếp tam giác
            \item Đường tròn nội tiếp tam giác
        \end{itemize}
        \end{minipage}
        &
        \begin{minipage}[t]{\linewidth}
        \begin{itemize}
            \item Nhận biết định nghĩa đường tròn ngoại tiếp tam giác.
            \item Xác định tâm và bán kính đường tròn ngoại tiếp tam giác, trong đó có tâm và bán kính đường tròn ngoại tiếp tam giác vuông, tam giác đều.
            \item Nhận biết định nghĩa đường tròn nội tiếp tam giác.
            \item Xác định tâm và bán kính đường tròn nội tiếp tam giác, trong đó có tâm và bán kính đường tròn nội tiếp tam giác đều.
        \end{itemize}
        \end{minipage}
    \end{tblr}
    \end{center}

    Cho trước một tam giác $ABC$. Bằng thước kẻ và compa, em có thể vẽ được một đường tròn đi qua ba đỉnh của tam giác và đường tròn tiếp xúc với cả ba cạnh của tam giác hay không?

    \subsubsection{Đường tròn ngoại tiếp một tam giác}

    \textbf{Khái niệm đường tròn ngoại tiếp tam giác}

    \textbf{HĐ1} Cho $d$ là đường trung trực của đoạn thẳng $AB$ và $O$ là một điểm trên $d$ (H.9.12). Hỏi đường tròn tâm $O$ đi qua $A$ thì có đi qua $B$ không?

    \begin{center}
    \begin{tikzpicture}
        \coordinate (A) at (-1.35,1.15);
        \coordinate (B) at (1.35,1.15);
        \coordinate (M) at ($(A)!0.5!(B)$);
        \coordinate (O) at (0,-0.75);
        \coordinate (Dtop) at (0,1.75);
        \coordinate (Dbot) at (0,-1.05);
        \draw (A)--(B);
        \draw (Dtop)--(Dbot);
        \draw[dashed] (A)--(O)--(B);
        \draw ($(A)!0.5!(M)+(0,-0.08)$)--($(A)!0.5!(M)+(0,0.08)$);
        \draw ($(M)!0.5!(B)+(0,-0.08)$)--($(M)!0.5!(B)+(0,0.08)$);
        \fill (O) circle (1.2pt);
        \node[above left=1pt of A] {$A$};
        \node[above right=1pt of B] {$B$};
        \node[below right=1pt of O] {$O$};
        \node[above right=1pt of Dtop] {$d$};
    \end{tikzpicture}

    \textit{Hình 9.12}
    \end{center}

    \answerline

    \textbf{HĐ2} Cho tam giác $ABC$ có ba đường trung trực đồng quy tại $O$ (H.9.13). Hãy giải thích tại sao đường tròn $(O;OA)$ đi qua ba đỉnh của tam giác $ABC$.

    \begin{center}
    \begin{tikzpicture}
        \coordinate (O) at (0,0);
        \coordinate (A) at (100:1.65);
        \coordinate (B) at (218:1.65);
        \coordinate (C) at (326:1.65);
        \coordinate (MAB) at ($(A)!0.5!(B)$);
        \coordinate (MBC) at ($(B)!0.5!(C)$);
        \coordinate (MCA) at ($(C)!0.5!(A)$);
        \draw (O) circle (1.65);
        \draw (A)--(B)--(C)--cycle;
        \draw[dashed] ($(O)!-1.15!(MAB)$)--($(O)!2.25!(MAB)$);
        \draw[dashed] ($(O)!-1.15!(MBC)$)--($(O)!2.25!(MBC)$);
        \draw[dashed] ($(O)!-1.15!(MCA)$)--($(O)!2.25!(MCA)$);
        \fill (O) circle (1.2pt);
        \node[above=1pt of A] {$A$};
        \node[below left=1pt of B] {$B$};
        \node[below right=1pt of C] {$C$};
        \node[below right=1pt of O] {$O$};
    \end{tikzpicture}

    \textit{Hình 9.13}
    \end{center}

    \answerline
    \answerline

    Đường tròn $(O;OA)$ như trên gọi là đường tròn ngoại tiếp tam giác $ABC$.

    Tổng quát, ta có định nghĩa sau:

    \begin{keyconcept}
        Đường tròn ngoại tiếp một tam giác là đường tròn đi qua ba đỉnh của tam giác đó.
    \end{keyconcept}

    Trong Hình 9.13, đường tròn $(O)$ ngoại tiếp tam giác $ABC$. Ta cũng nói tam giác $ABC$ nội tiếp đường tròn $(O)$, hay $(O)$ là đường tròn ngoại tiếp tam giác $ABC$. Tâm $O$ là giao điểm của ba đường trung trực của tam giác $ABC$.

    \textbf{Hãy kể tên bốn tam giác nội tiếp đường tròn $(O)$ trong Hình 9.14.}

    \begin{center}
    \begin{tikzpicture}
        \coordinate (O) at (0,0);
        \coordinate (B) at (132:1.5);
        \coordinate (C) at (222:1.5);
        \coordinate (M) at (25:1.5);
        \coordinate (N) at (-18:1.5);
        \coordinate (A) at (3.15,0.15);
        \draw (O) circle (1.5);
        \draw (B)--(C)--(M)--(B);
        \draw (B)--(N)--(C);
        \draw (M)--(N);
        \draw (B)--(A);
        \draw (C)--(A);
        \fill (O) circle (1.2pt);
        \node[above left=1pt of B] {$B$};
        \node[below left=1pt of C] {$C$};
        \node[above right=1pt of M] {$M$};
        \node[below right=1pt of N] {$N$};
        \node[right=1pt of A] {$A$};
        \node[above left=1pt of O] {$O$};
    \end{tikzpicture}

    \textit{Hình 9.14}
    \end{center}

    \answerline

    \textbf{Đường tròn ngoại tiếp tam giác vuông}

    \textbf{HĐ3} Cho tam giác $ABC$ vuông tại đỉnh $A$. Gọi $N$, $P$ lần lượt là trung điểm của các cạnh $AB$ và $AC$ (H.9.15).
    \begin{enumerate}[label=\alph*)]
        \item Vẽ hai đường trung trực $a$, $b$ của các cạnh $AB$, $AC$, cắt nhau tại $M$.

        \answerline
        \item Hãy giải thích vì sao $MN$, $MP$ là các đường trung bình của tam giác $ABC$.

        \answerline
        \answerline
        \item Hãy giải thích vì sao $M$ là trung điểm của $BC$, từ đó suy ra đường tròn ngoại tiếp của tam giác $ABC$ có tâm $M$ và bán kính $MB=MC=\dfrac{BC}{2}$.

        \answerline
        \answerline
    \end{enumerate}

    \begin{center}
    \begin{tikzpicture}
        \coordinate (A) at (0,0);
        \coordinate (B) at (0,3);
        \coordinate (C) at (4,0);
        \coordinate (M) at ($(B)!0.5!(C)$);
        \coordinate (N) at ($(A)!0.5!(B)$);
        \coordinate (P) at ($(A)!0.5!(C)$);
        \draw (M) circle ({sqrt(6.25)});
        \draw (A)--(B)--(C)--cycle;
        \draw (-0.65,1.5)--(2.7,1.5);
        \draw (2,-0.65)--(2,2.15);
        \pic[draw, angle radius=3.5mm] {right angle=B--A--C};
        \pic[draw, angle radius=2.8mm] {right angle=A--N--M};
        \pic[draw, angle radius=2.8mm] {right angle=A--P--M};
        \node[below left=1pt of A] {$A$};
        \node[above left=1pt of B] {$B$};
        \node[below right=1pt of C] {$C$};
        \node[above right=1pt of M] {$M$};
        \node[left=1pt of N] {$N$};
        \node[below=1pt of P] {$P$};
        \node[left=1pt] at (-0.65,1.5) {$a$};
        \node[below=1pt] at (2,-0.65) {$b$};
    \end{tikzpicture}

    \textit{Hình 9.15}
    \end{center}

    \begin{keyconcept}
        Đường tròn ngoại tiếp tam giác vuông có tâm là trung điểm của cạnh huyền và bán kính bằng một nửa cạnh huyền.
    \end{keyconcept}

    \textbf{Ví dụ 1} Cho tam giác $ABC$ vuông tại $A$ có $AB=2$ cm, $AC=4$ cm. Vẽ đường tròn $(O;R)$ ngoại tiếp tam giác $ABC$ và tính bán kính $R$.

    \textbf{Giải} (H.9.16)

    Lấy $O$ là trung điểm của $BC$ và vẽ đường tròn $(O)$ đi qua $A$. Khi đó, $(O)$ là đường tròn ngoại tiếp tam giác $ABC$.

    Theo định lí Pythagore, ta có:
    \[
        BC^2=AB^2+AC^2=4+16=20\text{ nên }BC=2\sqrt5\text{ (cm)}.
    \]
    Vậy đường tròn $(O)$ có bán kính $R=\dfrac{BC}{2}=\sqrt5$ cm.

    \begin{center}
    \begin{tikzpicture}
        \coordinate (A) at (0,0);
        \coordinate (B) at (0,2);
        \coordinate (C) at (4,0);
        \coordinate (O) at ($(B)!0.5!(C)$);
        \draw (O) circle ({sqrt(5)});
        \draw (A)--(B)--(C)--cycle;
        \pic[draw, angle radius=3.2mm] {right angle=B--A--C};
        \fill (O) circle (1.2pt);
        \node[below left=1pt of A] {$A$};
        \node[above left=1pt of B] {$B$};
        \node[below right=1pt of C] {$C$};
        \node[above right=1pt of O] {$O$};
    \end{tikzpicture}

    \textit{Hình 9.16}
    \end{center}

    \textbf{Luyện tập 1} Cho tam giác $ABC$ có $AC=3$ cm, $AB=4$ cm và $BC=5$ cm. Tính bán kính của đường tròn ngoại tiếp tam giác $ABC$.

    \answerline
    \answerline

    \textbf{Đường tròn ngoại tiếp tam giác đều}

    \textbf{HĐ4}
    \begin{enumerate}[label=\alph*)]
        \item Vẽ tam giác đều $ABC$. Hãy trình bày cách xác định tâm của đường tròn ngoại tiếp tam giác $ABC$ và vẽ đường tròn đó.

        \answerline
        \answerline
        \item Giải thích vì sao tâm $O$ của đường tròn ngoại tiếp tam giác $ABC$ trùng với trọng tâm của tam giác đó (H.9.17).

        \answerline
        \answerline
        \item Giải thích vì sao $\widehat{OBM}=30^\circ$ và $OB=\dfrac{\sqrt3}{3}BC$ (với $M$ là trung điểm của $BC$) (H.9.17).

        \answerline
        \answerline
    \end{enumerate}

    \begin{center}
    \begin{tikzpicture}
        \coordinate (O) at (0,0);
        \coordinate (A) at (90:1.75);
        \coordinate (B) at (210:1.75);
        \coordinate (C) at (330:1.75);
        \coordinate (M) at ($(B)!0.5!(C)$);
        \coordinate (N) at ($(A)!0.5!(C)$);
        \coordinate (P) at ($(A)!0.5!(B)$);
        \draw (O) circle (1.75);
        \draw (A)--(B)--(C)--cycle;
        \draw (A)--(M) (B)--(N) (C)--(P);
        \fill (O) circle (1.2pt);
        \node[above=1pt of A] {$A$};
        \node[below left=1pt of B] {$B$};
        \node[below right=1pt of C] {$C$};
        \node[above right=1pt of O] {$O$};
        \node[below=1pt of M] {$M$};
    \end{tikzpicture}

    \textit{Hình 9.17}
    \end{center}

    \begin{keyconcept}
        Đường tròn ngoại tiếp tam giác đều cạnh $a$ có tâm là trọng tâm của tam giác đó và bán kính bằng $\dfrac{\sqrt3}{3}a$.
    \end{keyconcept}

    \textbf{Ví dụ 2} Cho tam giác đều $ABC$ có độ dài cạnh bằng $3$ cm. Vẽ đường tròn $(O;R)$ ngoại tiếp tam giác $ABC$. Tính bán kính $R$.

    \textbf{Giải} (H.9.18)

    Lấy $O$ là giao điểm của ba đường trung tuyến của tam giác $ABC$ và vẽ đường tròn $(O)$ đi qua $A$. Đường tròn $(O)$ là đường tròn ngoại tiếp tam giác $ABC$ có bán kính
    \[
        R=\dfrac{\sqrt3}{3}BC=\sqrt3\text{ cm}.
    \]

    \begin{center}
    \begin{tikzpicture}
        \coordinate (O) at (0,0);
        \coordinate (A) at (90:1.75);
        \coordinate (B) at (210:1.75);
        \coordinate (C) at (330:1.75);
        \coordinate (M) at ($(B)!0.5!(C)$);
        \coordinate (N) at ($(A)!0.5!(C)$);
        \coordinate (P) at ($(A)!0.5!(B)$);
        \draw (O) circle (1.75);
        \draw (A)--(B)--(C)--cycle;
        \draw (A)--(M) (B)--(N) (C)--(P);
        \fill (O) circle (1.2pt);
        \node[above=1pt of A] {$A$};
        \node[below left=1pt of B] {$B$};
        \node[below right=1pt of C] {$C$};
        \node[above right=1pt of O] {$O$};
    \end{tikzpicture}

    \textit{Hình 9.18}
    \end{center}

    \textbf{Luyện tập 2} Cho tam giác đều $ABC$ nội tiếp đường tròn $(O)$ có bán kính bằng $4$ cm. Tính độ dài các cạnh của tam giác.

    \answerline
    \answerline

    \subsubsection{Đường tròn nội tiếp một tam giác}

    \textbf{Đường tròn nội tiếp tam giác}

    \textbf{HĐ5} Cho tam giác $ABC$ có ba đường phân giác đồng quy tại điểm $I$. Gọi $D$, $E$, $F$ lần lượt là chân các đường vuông góc kẻ từ $I$ xuống các cạnh $BC$, $CA$ và $AB$ (H.9.19).
    \begin{enumerate}[label=\alph*)]
        \item Hãy giải thích vì sao các điểm $D$, $E$, $F$ cùng nằm trên một đường tròn có tâm $I$.

        \answerline
        \answerline
        \item Gọi $(I)$ là đường tròn trên. Hãy giải thích vì sao $(I)$ tiếp xúc với các cạnh của tam giác $ABC$.

        \answerline
        \answerline
    \end{enumerate}

    \begin{center}
    \begin{tikzpicture}
        \coordinate (A) at (0,3);
        \coordinate (B) at (-2,0);
        \coordinate (C) at (2,0);
        \pgfmathsetmacro{\rin}{6/(2+sqrt(13))}
        \coordinate (I) at (0,\rin);
        \coordinate (D) at (0,0);
        \coordinate (E) at ($(A)!(I)!(C)$);
        \coordinate (F) at ($(A)!(I)!(B)$);
        \draw (A)--(B)--(C)--cycle;
        \draw (I) circle (\rin);
        \draw (A)--(I) (B)--(I) (C)--(I);
        \draw (I)--(D) (I)--(E) (I)--(F);
        \pic[draw, angle radius=2.5mm] {right angle=I--D--C};
        \pic[draw, angle radius=2.5mm] {right angle=I--E--C};
        \pic[draw, angle radius=2.5mm] {right angle=I--F--B};
        \fill (I) circle (1.2pt);
        \node[above=1pt of A] {$A$};
        \node[below left=1pt of B] {$B$};
        \node[below right=1pt of C] {$C$};
        \node[above right=1pt of I] {$I$};
        \node[below=1pt of D] {$D$};
        \node[right=1pt of E] {$E$};
        \node[left=1pt of F] {$F$};
    \end{tikzpicture}

    \textit{Hình 9.19}
    \end{center}

    Đường tròn $(I)$ như trên được gọi là đường tròn nội tiếp tam giác $ABC$ với các tiếp điểm trên các cạnh $BC$, $CA$, $AB$ lần lượt là $D$, $E$, $F$. Ta cũng nói tam giác $ABC$ ngoại tiếp đường tròn $(I)$ (H.9.19).

    Tổng quát, ta có định nghĩa sau:

    \begin{keyconcept}
        Đường tròn tiếp xúc với ba cạnh của tam giác được gọi là đường tròn nội tiếp tam giác. Tâm của đường tròn gọi là tâm nội tiếp đường tròn. Tâm đường tròn nội tiếp tam giác là giao điểm ba đường phân giác của tam giác.
    \end{keyconcept}

    Đường tròn tiếp xúc với một cạnh của tam giác nghĩa là tiếp xúc với đường thẳng chứa cạnh đó và có tiếp điểm nằm trên cạnh đó.

    \textbf{Mỗi tam giác có bao nhiêu đường tròn nội tiếp? Có bao nhiêu tam giác cùng ngoại tiếp một đường tròn?}

    \answerline
    \answerline

    \textbf{Đường tròn nội tiếp tam giác đều}

    \textbf{HĐ6} Cho tam giác đều $ABC$ có trọng tâm $G$.
    \begin{enumerate}[label=\alph*)]
        \item Giải thích vì sao $G$ cũng là tâm đường tròn nội tiếp tam giác $ABC$.

        \answerline
        \answerline
        \item Từ đó, giải thích vì sao bán kính đường tròn nội tiếp tam giác $ABC$ bằng một nửa bán kính đường tròn ngoại tiếp tam giác $ABC$ và bằng $\dfrac{\sqrt3}{6}BC$.

        \answerline
        \answerline
    \end{enumerate}

    \begin{keyconcept}
        Đường tròn nội tiếp tam giác đều cạnh $a$ có tâm là trọng tâm của tam giác đó và bán kính bằng $\dfrac{\sqrt3}{6}a$.
    \end{keyconcept}

    \textbf{Ví dụ 3} Cho tam giác $ABC$ ngoại tiếp đường tròn $(I)$. Biết rằng $\widehat{A}=40^\circ$, $\widehat{B}=60^\circ$. Tính số đo của các góc $BIC$, $CIA$ và $AIB$.

    \textbf{Giải} (H.9.20)

    Vì tổng ba góc của tam giác $ABC$ bằng $180^\circ$ nên
    \[
        \widehat{ACB}=180^\circ-\widehat{BAC}-\widehat{ABC}=80^\circ.
    \]

    Vì tam giác $ABC$ ngoại tiếp đường tròn $(I)$ nên $I$ là giao điểm các đường phân giác của tam giác $ABC$. Do đó, ta có: $\widehat{CAI}=\widehat{BAI}=\dfrac{\widehat{BAC}}{2}=20^\circ$; $\widehat{CBI}=\widehat{ABI}=\dfrac{\widehat{ABC}}{2}=30^\circ$; $\widehat{ACI}=\widehat{BCI}=\dfrac{\widehat{ACB}}{2}=40^\circ$.

    Vì tổng các góc trong tam giác $BIC$ bằng $180^\circ$ nên
    \[
        \widehat{BIC}=180^\circ-\widehat{CBI}-\widehat{BCI}=180^\circ-30^\circ-40^\circ=110^\circ.
    \]

    Tương tự, $\widehat{CIA}=180^\circ-\widehat{ACI}-\widehat{CAI}=120^\circ$ và $\widehat{AIB}=180^\circ-\widehat{ABI}-\widehat{BAI}=130^\circ$.

    \begin{center}
    \begin{tikzpicture}
        \coordinate (B) at (0,0);
        \coordinate (C) at (4,0);
        \path[name path=BA] (B)--++(60:5);
        \path[name path=CA] (C)--++(100:5);
        \path[name intersections={of=BA and CA, by=A}];
        \path[name path=BI] (B)--++(30:5);
        \path[name path=CI] (C)--++(140:5);
        \path[name intersections={of=BI and CI, by=I}];
        \pgfextracty{\dimen0}{\pgfpointanchor{I}{center}}
        \pgfmathsetlengthmacro{\rinex}{\dimen0}
        \draw (A)--(B)--(C)--cycle;
        \draw (A)--(I) (B)--(I) (C)--(I);
        \draw (I) circle (\rinex);
        \fill (I) circle (1.2pt);
        \node[above=1pt of A] {$C$};
        \node[below left=1pt of B] {$A$};
        \node[below right=1pt of C] {$B$};
        \node[below=1pt of I] {$I$};
    \end{tikzpicture}

    \textit{Hình 9.20}
    \end{center}

    \textbf{Thực hành} Vẽ đường tròn nội tiếp của tam giác $ABC$ bằng thước kẻ và compa theo các bước sau:
    \begin{itemize}
        \item Vẽ tia phân giác góc $B$ như sau: Dùng compa vẽ một cung tròn tâm $B$ cắt hai cạnh $BC$, $BA$ lần lượt tại $X$ và $Y$. Vẽ hai cung tròn tâm $X$, $Y$ có cùng bán kính, hai cung này cắt nhau tại một điểm $Z$ khác $B$. Kẻ tia $BZ$ ta được tia phân giác góc $B$.
        \item Tương tự, vẽ tia phân giác góc $C$, cắt tia $BZ$ tại $I$.
        \item Vẽ đường cao $ID$ từ $I$ xuống $BC$ ($D$ thuộc $BC$). Vẽ đường tròn $(I;ID)$ (H.9.21).
    \end{itemize}

    \begin{center}
    \begin{tikzpicture}
        \coordinate (A) at (0,3.2);
        \coordinate (B) at (-2.1,0);
        \coordinate (C) at (2.4,0);
        \path[name path=BI] (B)--($(B)!5cm!0.5:(A)$);
        % Hai tia phân giác được dựng bằng giao điểm của các đường phân giác trong.
        \coordinate (X) at ($(B)!0.34!(C)$);
        \coordinate (Y) at ($(B)!0.34!(A)$);
        \coordinate (XC) at ($(C)!0.34!(B)$);
        \coordinate (YC) at ($(C)!0.34!(A)$);
        \path let \p1=($(X)-(Y)$) in \pgfextra{\pgfmathsetmacro{\rB}{veclen(\x1,\y1)}};
        \path let \p1=($(XC)-(YC)$) in \pgfextra{\pgfmathsetmacro{\rC}{veclen(\x1,\y1)}};
        \path[name path=arcBX] (X) circle ({1.05*\rB});
        \path[name path=arcBY] (Y) circle ({1.05*\rB});
        \path[name intersections={of=arcBX and arcBY, by={Z,Zout}}];
        \path[name path=bisB] (B)--(Z);
        \path[name path=arcCX] (XC) circle ({1.05*\rC});
        \path[name path=arcCY] (YC) circle ({1.05*\rC});
        \path[name intersections={of=arcCX and arcCY, by={W,Wout}}];
        \path[name path=bisC] (C)--(W);
        \path[name intersections={of=bisB and bisC, by=I}];
        \coordinate (D) at ($(B)!(I)!(C)$);
        \coordinate (E) at ($(A)!(I)!(C)$);
        \draw (A)--(B)--(C)--cycle;
        \draw (B)--(I) (C)--(I);
        \draw (I)--(D);
        \draw (I) circle ({veclen(0,0)});
        \path let \p1=($(I)-(D)$) in \pgfextra{\pgfmathsetmacro{\rI}{veclen(\x1,\y1)}};
        \draw (I) circle (\rI);
        \draw[dashed] (B) ++(0:0.72) arc[start angle=0,end angle=57,radius=0.72];
        \draw[dashed] (X) circle ({1.05*\rB});
        \draw[dashed] (Y) circle ({1.05*\rB});
        \draw[dashed] (XC) circle ({1.05*\rC});
        \draw[dashed] (YC) circle ({1.05*\rC});
        \pic[draw, angle radius=2.8mm] {right angle=I--D--C};
        \node[above=1pt of A] {$A$};
        \node[below left=1pt of B] {$B$};
        \node[below right=1pt of C] {$C$};
        \node[below=1pt of D] {$D$};
        \node[below=1pt of X] {$X$};
        \node[left=1pt of Y] {$Y$};
        \node[above left=1pt of Z] {$Z$};
        \node[above right=1pt of I] {$I$};
    \end{tikzpicture}

    \textit{Hình 9.21}
    \end{center}

    Khi đó đường tròn $(I;ID)$ là đường tròn nội tiếp tam giác $ABC$ cần vẽ.

    \textbf{Luyện tập 3} Cho tam giác đều $ABC$ (H.9.22).
    \begin{enumerate}[label=\alph*)]
        \item Vẽ đường tròn $(I;r)$ nội tiếp tam giác $ABC$.

        \answerline
        \item Biết rằng $BC=4$ cm, hãy tính bán kính $r$.

        \answerline
        \answerline
    \end{enumerate}

    \begin{center}
    \begin{tikzpicture}
        \coordinate (B) at (-2,0);
        \coordinate (C) at (2,0);
        \coordinate (A) at (0,{2*sqrt(3)});
        \coordinate (I) at (0,{2*sqrt(3)/3});
        \pgfmathsetmacro{\req}{2*sqrt(3)/3}
        \draw (A)--(B)--(C)--cycle;
        \draw (I) circle (\req);
        \draw (A)--(I) (B)--(I) (C)--(I);
        \node[above=1pt of A] {$A$};
        \node[below left=1pt of B] {$B$};
        \node[below right=1pt of C] {$C$};
        \node[below=1pt of I] {$I$};
        \node[below=2pt] at ($(B)!0.5!(C)$) {$4$ cm};
    \end{tikzpicture}

    \textit{Hình 9.22}
    \end{center}

    \textbf{Em có biết?} \textit{(Đọc thêm)}

    \textbf{ĐƯỜNG TRÒN BÀNG TIẾP TAM GIÁC}

    Chúng ta biết rằng đường tròn nội tiếp một tam giác thì tiếp xúc với ba cạnh của tam giác và nằm bên trong tam giác. Ngoài ra, mỗi tam giác cho trước còn có ba đường tròn khác cũng tiếp xúc với ba đường thẳng chứa các cạnh của tam giác, nhưng chúng nằm ngoài tam giác. Các đường tròn đó gọi là các đường tròn bàng tiếp.

    Cụ thể, ta có định nghĩa sau: Đường tròn bàng tiếp của tam giác $ABC$ trong góc $A$ là đường tròn tiếp xúc với cạnh $BC$ và tiếp xúc với các tia đối của tia $BA$ và tia $CA$.

    Tương tự, ta cũng có đường tròn bàng tiếp trong các góc $B$ và $C$ của tam giác $ABC$.

    \textit{Vẽ đường tròn bàng tiếp.} Ta vẽ đường tròn bàng tiếp trong góc $A$ của tam giác $ABC$ như sau:
    \begin{itemize}
        \item Vẽ hai đường phân giác ngoài góc $B$ và $C$ của tam giác $ABC$. Gọi $I_a$ là giao điểm của hai đường phân giác ngoài này. Khi đó, $I_a$ nằm ngoài tam giác $ABC$ và cách đều các đường thẳng $AB$, $BC$, $CA$.
        \item Kẻ đường thẳng qua $I_a$ vuông góc với $AC$ và cắt $AC$ tại điểm $E$. Vẽ đường tròn $(I_a;I_aE)$. Khi đó, đường tròn $(I_a;I_aE)$ sẽ tiếp xúc với cạnh $BC$ và tiếp xúc với các tia đối của tia $BA$, $CA$.
    \end{itemize}

    Vậy $(I_a;I_aE)$ là đường tròn bàng tiếp trong góc $A$ của tam giác $ABC$ như hình bên.

    \begin{center}
    \begin{tikzpicture}
        \coordinate (A) at (0,3);
        \coordinate (B) at (-2,0);
        \coordinate (C) at (2,0);
        \pgfmathsetmacro{\rex}{6/(sqrt(13)-2)}
        \coordinate (Ia) at (0,-\rex);
        \coordinate (D) at (0,0);
        \coordinate (E) at ($(A)!(Ia)!(C)$);
        \draw (A)--(B)--(C)--cycle;
        \draw (B)--(Ia)--(C);
        \draw (Ia) circle (\rex);
        \draw (Ia)--(E);
        \pic[draw, angle radius=2.8mm] {right angle=Ia--E--C};
        \fill (Ia) circle (1.2pt);
        \node[above=1pt of A] {$A$};
        \node[above left=1pt of B] {$B$};
        \node[above right=1pt of C] {$C$};
        \node[below=1pt of Ia] {$I_a$};
        \node[right=1pt of E] {$E$};
    \end{tikzpicture}
    \end{center}
\end{lesson}

\begin{exercises}
    \subsection{Bài tập}

    \begin{questions}[resume=SGK]
        \item Cho đường tròn $(O)$ ngoại tiếp tam giác $ABC$. Tính bán kính của $(O)$, biết rằng $ABC$ là tam giác vuông cân tại $A$ và có độ dài cạnh bên bằng $2\sqrt2$ cm.

        \answerline
        \answerline

        \item Cho tam giác đều $ABC$ nội tiếp đường tròn $(O)$. Biết rằng đường tròn $(O)$ có bán kính bằng $3$ cm. Tính diện tích tam giác $ABC$.

        \answerline
        \answerline

        \item Cho tam giác $ABC$ nội tiếp đường tròn $(O)$. Gọi $H$ là trực tâm của tam giác $ABC$. Chứng minh rằng $\widehat{BAH}=\widehat{OAC}$.

        \answerline
        \answerline
        \answerline

        \item Cho đường tròn $(I)$ nội tiếp tam giác $ABC$ với các tiếp điểm trên các cạnh $AB$, $AC$ lần lượt là $E$, $F$. Chứng minh rằng $\widehat{EIF}+\widehat{BAC}=180^\circ$.

        \answerline
        \answerline
        \answerline

        \item Cho tam giác đều $ABC$ ngoại tiếp đường tròn $(I)$. Tính độ dài các cạnh của tam giác $ABC$ biết rằng bán kính của $(I)$ bằng $1$ cm.

        \answerline
        \answerline

        \item Người ta muốn làm một khung gỗ hình tam giác đều để đặt vừa khít một chiếc đồng hồ hình tròn có đường kính $30$ cm (H.9.23). Hỏi độ dài các cạnh (phía bên trong) của khung gỗ phải bằng bao nhiêu?

        % TODO(HINH-NGUON): Hình 9.23 là ảnh/minh hoạ chiếc đồng hồ trong khung gỗ ở SGK trang 76; bổ sung asset/crop đúng từ nguồn, không redraw xấp xỉ bằng TikZ.

        \answerline
        \answerline
    \end{questions}
\end{exercises}

\begin{practice}
    \subsection{Luyện tập}

    \begin{questions}[resume=SBT]
        \item Cho tam giác $ABC$ nội tiếp đường tròn $(O)$. Biết rằng $\widehat{ACB}=50^\circ$, $\widehat{ABC}=70^\circ$, tính số đo các cung nhỏ $\widehat{BC}$, $\widehat{CA}$, $\widehat{AB}$ của đường tròn $(O)$.

        \answerline
        \answerline

        \item Cho tam giác $ABC$ cân tại $A$ và nội tiếp đường tròn $(O)$. Tính số đo các góc của tam giác $ABC$, biết rằng $\widehat{BOC}=100^\circ$.

        \answerline
        \answerline

        \item Tính bán kính và chu vi của đường tròn ngoại tiếp tam giác $ABC$ có ba cạnh $AB=6$ cm, $AC=8$ cm và $BC=10$ cm.

        \answerline
        \answerline

        \item Tính chu vi và diện tích của tam giác đều nội tiếp một đường tròn bán kính $3$ cm.

        \answerline
        \answerline

        \item Tính chu vi và diện tích của tam giác đều ngoại tiếp một đường tròn bán kính $3$ cm.

        \answerline
        \answerline

        \item Cho tam giác $ABC$ vuông tại $A$ có $AB=4$ cm, $AC=6$ cm. Tính bán kính của đường tròn nội tiếp tam giác $ABC$.

        \answerline
        \answerline

        \item Cho tam giác $ABC$ vuông tại $A$ nội tiếp đường tròn $(O)$ có bán kính $2$ cm. Biết rằng $AC=2$ cm, tính số đo các góc của tam giác $ABC$.

        \answerline
        \answerline

        \item Cho tam giác $ABC$ có trực tâm $H$ và nội tiếp đường tròn $(O)$. Chứng minh rằng:
        \begin{questions}
            \item $\widehat{OBC}=90^\circ-\widehat{BAC}$;

            \answerline
            \answerline
            \item $\widehat{BAH}=\widehat{OAC}$.

            \answerline
            \answerline
        \end{questions}

        \item Cho tam giác $ABC$ ngoại tiếp đường tròn $(I)$. Chứng minh rằng:
        \[
            \widehat{BIC}=90^\circ+\dfrac{\widehat{BAC}}{2};\qquad
            \widehat{CIA}=90^\circ+\dfrac{\widehat{CBA}}{2};\qquad
            \widehat{AIB}=90^\circ+\dfrac{\widehat{ACB}}{2}.
        \]

        \answerline
        \answerline
        \answerline

        \item Cho tam giác nhọn $ABC$ nội tiếp đường tròn $(O)$. Cho điểm $M$ trên cạnh $BC$ của tam giác $ABC$ và điểm $D$ trên cung nhỏ $BC$ của $(O)$ sao cho $\widehat{BAD}=\widehat{MAC}$. Chứng minh rằng $\triangle AMB\sim\triangle ACD$.

        \answerline
        \answerline
        \answerline

        \item Cho tam giác $ABC$ vuông tại $B$ nội tiếp đường tròn $(O)$ và đường kính $BD$. Tính số đo của góc $BAC$, biết rằng $\widehat{BAC}=2\widehat{CBD}$.

        \answerline
        \answerline

        \item Cho tam giác $ABC$ nội tiếp đường tròn $(O)$ và ngoại tiếp đường tròn $(I)$. Tia $AI$ cắt $(O)$ tại $X$ (khác $A$). Chứng minh rằng $X$ là tâm của đường tròn ngoại tiếp tam giác $BIC$.

        \answerline
        \answerline
        \answerline

        \item Cho $\triangle ABC\sim\triangle A'B'C'$ với tỉ số đồng dạng $k>0$. Gọi $(O,R)$ và $(O',R')$ lần lượt là đường tròn ngoại tiếp các tam giác $ABC$ và $A'B'C'$. Gọi $(I,r)$ và $(I',r')$ lần lượt là đường tròn nội tiếp các tam giác $ABC$ và $A'B'C'$. Chứng minh rằng
        \[
            \dfrac{R}{R'}=\dfrac{r}{r'}=k.
        \]

        \answerline
        \answerline
        \answerline
    \end{questions}
\end{practice}
